\documentclass[11pt]{amsart}
\usepackage[T1]{fontenc}
\usepackage{lmodern,microtype,amsmath,amssymb,amsthm,mathtools,enumitem,booktabs}
\usepackage[margin=1in]{geometry}
\usepackage[hidelinks]{hyperref}
\usepackage{xurl}
\newtheorem{theorem}{Theorem}[section]
\newtheorem{proposition}[theorem]{Proposition}
\newtheorem{lemma}[theorem]{Lemma}
\newtheorem{corollary}[theorem]{Corollary}
\theoremstyle{remark}
\newtheorem{remark}[theorem]{Remark}
\numberwithin{equation}{section}
\newcommand{\E}{\mathbb E}
\newcommand{\Pp}{\mathbb P}
\newcommand{\DTCS}{\mathrm{DTCS}}
\hypersetup{pdftitle={The drawn length of the critical beta-splitting tree grows linearly},pdfauthor={Yitzchak Shmalo}}
\title[The drawn length of the critical beta-splitting tree]{The drawn length of the critical\\ beta-splitting tree grows linearly}
\author{Yitzchak Shmalo}
\address{Einstein Institute of Mathematics, The Hebrew University of Jerusalem, Jerusalem, Israel}
\date{12 October 2026}
\subjclass[2020]{60C05, 05C05, 60J10, 92B10}
\keywords{Critical beta-splitting tree, cladogram, drawn length, harmonic descent chain, occupation measure}

\begin{document}
\begin{abstract}
In the critical beta-splitting model of a random binary tree with $n$ leaves, a clade of $m$ leaves splits into clades of $i$ and $m-i$ leaves with probability proportional to $1/(i(m-i))$. Aldous and Janson asked for a proof that the expected drawn length $L_n$ of the cladogram of this tree grows roughly like $n\log^4 n$. We show that $L_n/n$ converges to the finite constant $\lambda=\sum_{m\ge2}a(m)c_m/m$, where $a(m)=6h_{m-1}/(\pi^2(m-1))$ is the limiting probability that a uniformly chosen leaf lies in a clade of exactly $m$ leaves and $c_m$ is the expected contribution of the root of a tree with $m$ leaves to the drawn length. Numerically $\lambda\approx4.88$, and $4.80<\lambda<5.38$. The proof combines an exponential moment bound for the height with a bound, uniform in $n$, on the time the harmonic descent chain spends in a dyadic block. A deterministic comparison with the sum of the leaf depths gives the weaker bound $L_n=O(n\log^2n)$ directly. The prediction $n\log^4n$ comes from a heuristic in which the expected number of clades of size $m$ is taken to be $na(n,m)$ instead of $(n/m)a(n,m)$.
\end{abstract}
\maketitle

\section{Introduction}\label{sec:intro}

Let $h_k=\sum_{j=1}^k1/j$ denote the harmonic numbers. The critical beta-splitting tree $T_n$, denoted $\DTCS(n)$ in \cite{AJ23}, is the random rooted binary tree with $n$ leaves constructed as follows. The root clade is the set of all $n$ leaves. A clade of $m\ge2$ leaves is split into a left clade of $i$ leaves and a right clade of $m-i$ leaves with probability
\begin{equation}\label{eq:q}
 q(m,i)=\frac{m}{2h_{m-1}\,i(m-i)},\qquad 1\le i\le m-1,
\end{equation}
independently of all other splits, and a clade of one leaf is a leaf of the tree. This is the case $\beta=-1$ of the beta-splitting family of Aldous \cite{Ald96}, and its structure has been studied in detail in recent years \cite{AP25,AJ23,AJL24,BCM25,Ma26}.

For a finite rooted binary tree $T$ we write $h(T)$ for its height, the largest number of edges on a path from the root to a leaf, so that a single leaf has height $0$, and $T_v$ for the subtree consisting of a vertex $v$ and its descendants. In the usual drawing of $T$ as a cladogram, the leaves sit at level $0$, an internal vertex $v$ sits at level $h(T_v)$, and $v$ is joined to each of its two children by a vertical line. The \emph{drawn length} $\ell(T)$ is the total length of these vertical lines. If $T_L$ and $T_R$ are the subtrees of the two children of the root, then $h(T)=1+\max(h(T_L),h(T_R))$ and
\begin{equation}\label{eq:drawn}
 \ell(T)=\ell(T_L)+\ell(T_R)+2+|h(T_L)-h(T_R)|,
\end{equation}
with $\ell(T)=0$ for a single leaf; this is the recursion (73) of \cite{AJ23}. Put
\[
 L_n=\E\,\ell(T_n),\qquad H_n=\E\,h(T_n).
\]
In \cite[Section 6.1.1]{AJ23} Aldous and Janson give a heuristic argument that $L_n$ is of order $n\log^4n$ and pose the following problem.

\begin{quote}
\textbf{Open Problem 14 of \cite{AJ23}.} Prove that $L_n$ grows roughly like $n\log^4n$.
\end{quote}

The order is in fact $n$. For $m\ge2$ let
\begin{equation}\label{eq:cm}
 c_m=\E\bigl[2+|h(T_L)-h(T_R)|\bigr],
\end{equation}
where $T_L$ and $T_R$ are the subtrees of the two children of the root of $T_m$, and let
\begin{equation}\label{eq:am}
 a(m)=\frac{6h_{m-1}}{\pi^2(m-1)},\qquad m\ge2.
\end{equation}

\begin{theorem}\label{thm:main}
As $n\to\infty$,
\[
 \frac{L_n}{n}\longrightarrow\lambda:=\sum_{m=2}^\infty\frac{a(m)\,c_m}{m}.
\]
The series converges, and $2<\lambda<\infty$. Numerically $4.800<\lambda<5.372$, and $\lambda\approx4.88$.
\end{theorem}

The numerical statements are explained in Section~\ref{sec:numerics}. A weaker upper bound, which already rules out the order $n\log^4n$, follows from a deterministic comparison. Let $S(T)$ be the sum of the depths of the leaves of $T$ (the Sackin index of $T$).

\begin{proposition}\label{prop:sackin}
For every finite rooted binary tree $T$ with $n$ leaves, $2(n-1)\le\ell(T)\le S(T)$. Consequently
\[
 2(n-1)\le L_n\le n\,\E D_n=\Bigl(\frac3{\pi^2}+o(1)\Bigr)n\log^2n,
\]
where $D_n$ is the depth of a uniformly chosen leaf of $T_n$.
\end{proposition}

The first statement is part of Lemma~\ref{lem:det} below, $\E S(T_n)=n\E D_n$ holds by averaging over the leaves, and $\E D_n\sim(3/\pi^2)\log^2n$ is proved by Aldous and Pittel \cite{AP25} (see also \cite[Theorem 3]{AJ23}).

The heuristic of \cite[(75)]{AJ23} starts from the expected number of clades of size $m$ in $T_n$, which is $(n/m)a(n,m)$, where $a(n,m)$ is the probability that a uniformly chosen leaf lies in a clade of exactly $m$ leaves; this identity is stated in \cite[Section 4.8]{AJ23}, and we prove it in Lemma~\ref{lem:clades}. The first line of (75) uses $n\,a(n,m)$ instead. The two counts differ in total: $\sum_{m=2}^n(n/m)a(n,m)=n-1$ is the number of internal vertices of $T_n$, while $\sum_{m=2}^n n\,a(n,m)=n\E D_n$. With the factor $1/m$ in place, the same computation leads to the convergent series of Theorem~\ref{thm:main}.

The proof of Theorem~\ref{thm:main} rests on two elementary estimates. Lemma~\ref{lem:height} bounds $H_n\le4h_{n-1}\log n$ through an exponential moment of $D_n$, and Lemma~\ref{lem:block} bounds the expected number of visits of the harmonic descent chain to a block $\{b+1,\dots,2b\}$ by $2h_{2b-1}$, uniformly in the starting point. Together with the convergence $a(n,m)\to a(m)$ for each fixed $m$, which is Theorem~18 of \cite{AJ23}, they give the convergence of $L_n/n=\sum_ma(n,m)c_m/m$ by dominated convergence. The convergence of $\sum_ma(n,m)f(m)$ to $\sum_ma(m)f(m)$ for nonnegative $f$ with $f(m)=O(m^{-\varepsilon})$ is a case of Ansatz~20 of \cite{AJ23} stated by Brandenberger, Chin and Mossel \cite[Section 1.1]{BCM25}, who also obtain rates for $a(n,m)\to a(m)$. Since $c_m/m=O(m^{-1/2})$ by Lemma~\ref{lem:height}, Theorem~\ref{thm:main} also follows from their statement; the argument given here uses only the convergence $a(n,m)\to a(m)$. Ma \cite{Ma26} proved that $h(T_n)/\log^2n$ converges almost surely and in $L^p$ to an explicit constant $C_{\mathrm{ht}}\approx0.976$, which answers Open Problem~4 of \cite{AJ23}; the cruder Lemma~\ref{lem:height} is all that is needed here. Additive shape parameters with explicit tolls, such as the Sackin and Colless indices, are studied across the beta-splitting family in \cite{Son26}.

Theorem~\ref{thm:main} is a statement about the mean. Whether $\ell(T_n)/n$ converges to $\lambda$ in probability, and what the fluctuations are, we leave open, as well as Open Problem~15 of \cite{AJ23} on the profile of the drawn width.

Section~\ref{sec:identities} collects the identities, Section~\ref{sec:estimates} proves the two estimates, Section~\ref{sec:proof} proves Theorem~\ref{thm:main}, and Section~\ref{sec:numerics} describes the computations, including a Lean proof of Lemma~\ref{lem:det}.

\section{Identities}\label{sec:identities}

For a finite rooted binary tree $T$ let $n(T)$ be its number of leaves and $I(T)$ its set of internal vertices, so that $|I(T)|=n(T)-1$.

\begin{lemma}\label{lem:det}
Let $T$ be a finite rooted binary tree with $n$ leaves. Then
\begin{enumerate}[label=\textup{(\alph*)}]
\item $\ell(T)=\sum_{v\in I(T)}h(T_v)+h(T)$;
\item $S(T)=\sum_{v\in I(T)}n(T_v)$;
\item $h(T)\le n-1$;
\item $2(n-1)\le\ell(T)\le S(T)$.
\end{enumerate}
\end{lemma}

\begin{proof}
(a) The vertical line from an internal vertex $v$ to a child $u$ has length $h(T_v)-h(T_u)$. In the sum over all vertical lines, every internal vertex appears twice as a parent, every vertex other than the root appears once as a child, and leaves have height $0$. Hence $\ell(T)=2\sum_{v\in I(T)}h(T_v)-\sum_{v\in I(T)\setminus\{\mathrm{root}\}}h(T_v)$, which is (a).
(b) The depth of a leaf is the number of its internal ancestors, so each $v\in I(T)$ is counted once for each of the $n(T_v)$ leaves below it.
(c) By induction, $h(T)=1+\max(h(T_L),h(T_R))\le\max(n(T_L),n(T_R))\le n-1$.
(d) By \eqref{eq:drawn} every internal vertex contributes at least $2$ to $\ell(T)$, which gives the lower bound. By (a), (c) applied to each $T_v$, and (b),
\[
 \ell(T)\le\sum_{v\in I(T)}\bigl(n(T_v)-1\bigr)+(n-1)=S(T)-|I(T)|+(n-1)=S(T).\qedhere
\]
\end{proof}

Fix $n$ and let $V$ be a leaf of $T_n$ chosen uniformly at random. The sizes of the clades containing $V$, read from the root down, form a decreasing sequence $n=J_0>J_1>\dots>J_{D_n}=1$, where $D_n$ is the depth of $V$. Given $J_t=k\ge2$, the clade of size $k$ splits as in \eqref{eq:q} and $V$, being uniform in it, falls into a part of size $j$ with probability $j/k$; since $q(k,j)=q(k,k-j)$,
\begin{equation}\label{eq:hd}
 \Pp(J_{t+1}=j\mid J_t=k)=2q(k,j)\frac jk=\frac1{h_{k-1}(k-j)},\qquad 1\le j\le k-1.
\end{equation}
So $(J_t)$ is the harmonic descent chain of \cite{AJL24} started at $n$ and stopped at $1$ (see also \cite[Section 3.1]{AJ23}). For $1\le m\le n$ let $a(n,m)$ be the probability that this chain visits $m$, and put $a(n,m)=0$ for $m>n$. Then $\E D_n=\sum_{m=2}^na(n,m)$.

\begin{lemma}\label{lem:clades}
Let $N_n(m)$ be the number of internal vertices $v$ of $T_n$ with $n(T_v)=m$. Then $\E N_n(m)=(n/m)\,a(n,m)$ for $2\le m\le n$.
\end{lemma}

\begin{proof}
Two clades of the same size are either equal or disjoint, so given $T_n$ the probability that $V$ lies in a clade of size $m$ is $mN_n(m)/n$. Taking expectations gives $a(n,m)=m\E N_n(m)/n$.
\end{proof}

\begin{lemma}\label{lem:mean}
For $n\ge2$,
\begin{align}
 L_n&=\sum_{m=2}^n\E N_n(m)\,c_m=n\sum_{m=2}^n\frac{a(n,m)\,c_m}{m},\label{eq:Lc}\\
 L_n&=n\sum_{m=2}^n\frac{a(n,m)\,H_m}{m}+H_n.\label{eq:LH}
\end{align}
Moreover $c_m=2H_m-2\sum_{i=1}^{m-1}q(m,i)H_i$ and $2\le c_m\le1+H_m$ for $m\ge2$.
\end{lemma}

\begin{proof}
Unrolling \eqref{eq:drawn} gives $\ell(T)=\sum_{v\in I(T)}g(T_v)$ with $g(T')=2+|h(T'_L)-h(T'_R)|$. Conditionally on the splits outside it, the subtree below a clade of size $m$ is distributed as $T_m$. Hence $\E\sum_{v\in I(T_n)}g(T_v)=\sum_m\E N_n(m)\,\E g(T_m)=\sum_m\E N_n(m)c_m$, and Lemma~\ref{lem:clades} gives \eqref{eq:Lc}. The same argument applied to Lemma~\ref{lem:det}(a) gives \eqref{eq:LH}. For the formula for $c_m$ use $|x-y|=2\max(x,y)-x-y$, $\E\max(h(T_L),h(T_R))=H_m-1$, and $\E h(T_L)=\E h(T_R)=\sum_iq(m,i)H_i$, which holds because $q(m,i)=q(m,m-i)$. The lower bound $c_m\ge2$ is clear, and $|x-y|\le\max(x,y)$ for $x,y\ge0$ gives $c_m\le2+H_m-1$.
\end{proof}

Since $\sum_{m=2}^n\E N_n(m)=n-1$, Lemma~\ref{lem:clades} also gives
\begin{equation}\label{eq:mass}
 \sum_{m=2}^n\frac{a(n,m)}m=\frac{n-1}n.
\end{equation}

\section{Two estimates}\label{sec:estimates}

\begin{lemma}\label{lem:height}
For $n\ge2$, $H_n\le4h_{n-1}\log n$. Consequently $c_m\le1+4h_{m-1}\log m$ for $m\ge2$.
\end{lemma}

\begin{proof}
Let $\theta=1/(2h_{n-1})$. For $1\le k\le n$ let $F(k)=\E_k e^{\theta D}$, where $D$ is the number of steps the chain \eqref{eq:hd} started at $k$ needs to reach $1$. We show $F(k)\le k$ by induction. Clearly $F(1)=1$. Let $2\le k\le n$ and assume $F(j)\le j$ for $j<k$. Since
\[
 \sum_{j=1}^{k-1}\frac{j}{h_{k-1}(k-j)}=\sum_{r=1}^{k-1}\frac{k-r}{h_{k-1}\,r}=k-\frac{k-1}{h_{k-1}},
\]
we get $F(k)\le e^\theta\sum_{j<k}\Pp(J_1=j\mid J_0=k)\,j=e^\theta k(1-x_k)$ with $x_k=(k-1)/(kh_{k-1})$. Now $x_k\ge1/(2h_{k-1})\ge\theta$ and $x_k<1$, so $e^\theta\le1/(1-\theta)\le1/(1-x_k)$ and $F(k)\le k$.

Pointwise $e^{\theta h(T_n)}$ is at most the sum of $e^{\theta d(v)}$ over the leaves $v$, where $d(v)$ is the depth of $v$, and the expectation of this sum is $n\E e^{\theta D_n}=nF(n)\le n^2$. By Jensen's inequality $e^{\theta H_n}\le\E e^{\theta h(T_n)}\le n^2$, that is, $H_n\le2\log n/\theta=4h_{n-1}\log n$. The bound on $c_m$ follows from Lemma~\ref{lem:mean}.
\end{proof}

\begin{lemma}\label{lem:block}
For all $n\ge2$ and $b\ge1$,
\[
 \sum_{m=b+1}^{2b}a(n,m)\le2h_{2b-1}.
\]
\end{lemma}

\begin{proof}
The sum is the expected number of visits of the chain \eqref{eq:hd} started at $n$ to $B=\{b+1,\dots,2b\}$. The chain is strictly decreasing, so its visits to $B$ are consecutive. From any $k\in B$ it jumps to $\{1,\dots,b\}$ with probability
\[
 \sum_{j=1}^b\frac1{h_{k-1}(k-j)}\ge\frac{b}{h_{k-1}(k-1)}\ge\frac{b}{(2b-1)h_{2b-1}}\ge\frac1{2h_{2b-1}}.
\]
Hence the number of visits to $B$ is stochastically dominated by a geometric random variable with mean $2h_{2b-1}$.
\end{proof}

\section{Proof of Theorem~\ref{thm:main}}\label{sec:proof}

Let $f(m)=c_m/m$. For $b\ge1$ and $b<m\le2b$, Lemma~\ref{lem:height} gives $f(m)\le(1+4h_{2b-1}\log(2b))/b$, and with Lemma~\ref{lem:block}, for every $n$,
\begin{equation}\label{eq:blockbound}
 \sum_{m=b+1}^{2b}a(n,m)f(m)\le\beta_b:=\frac{2h_{2b-1}\bigl(1+4h_{2b-1}\log(2b)\bigr)}{b}.
\end{equation}
Since $h_{2^{k+1}-1}\le1+(k+1)\log2$, the numbers $\beta_{2^k}$ are of order $k^3 2^{-k}$, and $\varepsilon_K:=\sum_{k\ge K}\beta_{2^k}$ is finite and tends to $0$ as $K\to\infty$. The blocks $\{2^k+1,\dots,2^{k+1}\}$, $k\ge0$, partition $\{2,3,\dots\}$, so \eqref{eq:blockbound} gives
\begin{equation}\label{eq:tail}
 \sum_{m>2^K}a(n,m)f(m)\le\varepsilon_K\qquad\text{for all }n\text{ and }K.
\end{equation}
Each block sum in \eqref{eq:blockbound} has finitely many terms, so letting $n\to\infty$ and using $a(n,m)\to a(m)$ \cite[Theorem 18]{AJ23} gives $\sum_{b<m\le2b}a(m)f(m)\le\beta_b$, and hence $\sum_{m>2^K}a(m)f(m)\le\varepsilon_K$. In particular $\lambda=\sum_ma(m)f(m)<\infty$. By \eqref{eq:Lc} and \eqref{eq:tail},
\[
 \Bigl|\frac{L_n}n-\lambda\Bigr|\le\sum_{m=2}^{2^K}|a(n,m)-a(m)|\,f(m)+2\varepsilon_K.
\]
The finite sum tends to $0$ as $n\to\infty$, and then $\varepsilon_K\to0$ as $K\to\infty$. Hence $L_n/n\to\lambda$.

For the lower bound, \eqref{eq:Lc} and \eqref{eq:mass} give
\[
 \frac{L_n}n=\frac{2(n-1)}n+\sum_{m=2}^n\frac{a(n,m)(c_m-2)}m.
\]
Since $0\le(c_m-2)/m\le f(m)$, the argument above applies to the last sum, and therefore $\lambda=2+\sum_{m\ge2}a(m)(c_m-2)/m\ge2+a(4)(c_4-2)/4$. A tree with $4$ leaves splits into parts of sizes $1$ and $3$ with probability $8/11$, and then $|h(T_L)-h(T_R)|=2$; otherwise both parts have height $1$. So $c_4=2+16/11$ and $\lambda>2$. The numerical bounds are proved in Section~\ref{sec:numerics}.\qed

\begin{remark}\label{rem:H}
By \eqref{eq:LH}, Lemma~\ref{lem:height}, which gives $H_n/n\to0$, and the same argument with $H_m/m$ in place of $f(m)$, also
\[
 \lambda=\sum_{m=2}^\infty\frac{a(m)\,H_m}{m}.
\]
The two series for $\lambda$ have different terms; their sums agree because $c_m=2H_m-2\sum_iq(m,i)H_i$ and $\sum_{m>i}a(m)q(m,i)/m=a(i)/(2i)$, the limit form of the recursion $\E N_n(i)=\mathbf 1\{i=n\}+2\sum_{m>i}\E N_n(m)q(m,i)$.
\end{remark}

\section{Numerical values}\label{sec:numerics}

The script \texttt{code/check.py} in the repository of the paper computes the distribution of $h(T_m)$ for $m\le4000$ from the recursion
\[
 \Pp(h(T_m)\le x)=\sum_{i=1}^{m-1}q(m,i)\,\Pp(h(T_i)\le x-1)\,\Pp(h(T_{m-i})\le x-1),
\]
truncated at $x=400$, beyond which the remaining probability for $m=4000$ is below $10^{-12}$. From it, it computes $H_m$, then $c_m$ by Lemma~\ref{lem:mean}, and $L_m$ by the recursion $L_m=2\sum_iq(m,i)L_i+c_m$, which follows from \eqref{eq:drawn} and the branching property. For $n\in\{10,100,1000,2000,4000\}$ it also computes the expected clade counts from the recursion $\E N_n(m)=\mathbf 1\{m=n\}+2\sum_{k>m}\E N_n(k)q(k,m)$ and the probabilities $a(n,m)$ from the chain \eqref{eq:hd}, and checks Lemma~\ref{lem:clades}, which holds to a relative error below $3\cdot10^{-15}$, and \eqref{eq:Lc}: the two computations of $L_n/n$ agree to within $10^{-9}$. In exact rational arithmetic, the script computes the joint law of $(h(T_m),\ell(T_m))$ for $m\le8$ by conditioning on the first split and checks Lemma~\ref{lem:clades} for $n\le12$ and \eqref{eq:Lc} for $n\le8$; Lemma~\ref{lem:det} is checked on all $6918$ ordered binary trees with at most $10$ leaves. Lemma~\ref{lem:height} holds for $n\le4000$ with largest ratio $H_n/(4h_{n-1}\log n)$ equal to $0.361$, and Lemma~\ref{lem:block} holds for $n=4000$ and every $b\le2000$ with largest ratio $0.407$. Tables~\ref{tab:exact} and~\ref{tab:float} show some of the values.

\begin{table}[ht]
\centering
\begin{tabular}{rcccccc}
\toprule
$m$ & $2$ & $3$ & $4$ & $5$ & $6$ & $7$\\
\midrule
$H_m$ & $1$ & $2$ & $\tfrac{30}{11}$ & $\tfrac{189}{55}$ & $\tfrac{30093}{7535}$ & $\tfrac{239091}{52745}$\\[2pt]
$L_m$ & $2$ & $5$ & $\tfrac{90}{11}$ & $\tfrac{646}{55}$ & $\tfrac{2103}{137}$ & $\tfrac{201977}{10549}$\\
\bottomrule
\end{tabular}
\medskip
\caption{Exact mean height and mean drawn length of $T_m$.}\label{tab:exact}
\end{table}

\begin{table}[ht]
\centering
\begin{tabular}{rrr}
\toprule
$n$ & $H_n$ & $L_n/n$\\
\midrule
$10$ & $5.8603$ & $3.0988$\\
$100$ & $18.5362$ & $4.4286$\\
$1000$ & $39.1932$ & $4.7877$\\
$2000$ & $47.1006$ & $4.8239$\\
$4000$ & $55.8147$ & $4.8458$\\
\bottomrule
\end{tabular}
\medskip
\caption{Mean height and mean drawn length per leaf.}\label{tab:float}
\end{table}

\subsection*{The constant \texorpdfstring{$\lambda$}{lambda}}
All terms of the series for $\lambda$ are positive, so the partial sum $\sum_{m=2}^{4000}a(m)c_m/m=4.80059$ is a lower bound. For $m>4000$ we use $c_m\le1+4h_m\log m$ from Lemma~\ref{lem:height}; the resulting series is summed numerically up to $m=10^8$ and bounded by an integral beyond, and its value is less than $0.5712$. Hence $4.800<\lambda<5.372$. For $1000\le m\le4000$ the values $c_m$ are fitted by $0.5756\log^2m-0.5604\log m+4.7753$ with largest error $6\cdot10^{-4}$, and using this fit for $m>4000$ gives $\lambda\approx4.877$. The series of Remark~\ref{rem:H}, with $H_m$ fitted in the same way, gives $4.877$ as well, although its partial sum ($4.7686$) and its estimated remainder ($0.1085$) are different. The values of $L_n/n$ in Table~\ref{tab:float} increase slowly towards $\lambda$; by Lemma~\ref{lem:height} the remainder $\sum_{m>n}a(m)c_m/m$ is at most of order $(\log n)^3/n$.

\subsection*{Lean}
The file \texttt{lean/DrawnLength.lean} proves Lemma~\ref{lem:det} for all finite rooted binary trees in Lean~4, using the core library only. The theorems depend on the axioms \texttt{propext} and \texttt{Quot.sound}. The probabilistic statements are not formalized.

\section*{Acknowledgements and funding}
The research presented in this paper was supported by the European Research
Council (ERC) under the European Union's Horizon Europe research and
innovation programme (grant agreement No.~101041711), by the Simons
Foundation as part of the Collaboration on the Mathematical and Scientific
Foundations of Deep Learning, by Heights Labs, by Convex Nexus Capital, by
the Israel Science Foundation (grant number 2258/19), and by the Israel
Science Foundation (ISF Grant 4101/25).

\section*{Declaration of competing interest}
The author holds equity in Heights Labs and in Convex Nexus Capital and
has a management role at Convex Nexus Capital. These entities are included
in the funding declaration above.

\section*{Declaration of generative AI and AI-assisted technologies}
The paper comes out of Math Brain, an ongoing project on the use of AI systems for research mathematics. Open Problem~14 of \cite{AJ23} was one of the open problems in probability on which the project's agents worked in September 2026. On 6 September 2026 an OpenAI Codex agent wrote the comparison of Proposition~\ref{prop:sackin}, observed that the heuristic (75) of \cite{AJ23} omits the factor $1/m$, and gave an argument for the linear limit in the form of Remark~\ref{rem:H}, based on a height bound through halving steps and the block estimate of Lemma~\ref{lem:block}. On 8 September 2026 a Claude agent (Anthropic) wrote a Lean proof of the inequality $\ell(T)\le S(T)$ and estimated $L_n/n$ by simulation. On 10 October 2026 an OpenAI Codex agent checked the argument again. On 12 October 2026 Claude Opus~5.5 (Anthropic), working in Claude Code, gave the exponential moment proof of Lemma~\ref{lem:height}, wrote the proofs in their present form, searched the literature, and wrote the computer checks, the Lean file of Section~\ref{sec:numerics} and the manuscript. The author directed the work and is responsible for the content of the paper and for decisions concerning its dissemination.

\section*{Data availability}
No data were used. The code and the Lean file are available at \url{https://github.com/yspennstate/beta-splitting-drawn-length}.

\begin{thebibliography}{99}

\bibitem{Ald96}
D.~Aldous,
Probability distributions on cladograms,
in: \emph{Random Discrete Structures} (D.~Aldous and R.~Pemantle, eds.),
IMA Vol. Math. Appl. \textbf{76}, Springer, New York, 1996, pp.~1--18.

\bibitem{AJ23}
D.~J.~Aldous and S.~Janson,
The critical beta-splitting random tree II: overview and open problems,
arXiv:2303.02529v3, 2025.

\bibitem{AJL24}
D.~J.~Aldous, S.~Janson and X.~Li,
The harmonic descent chain,
\emph{Electron. Commun. Probab.} \textbf{29} (2024), 1--10.

\bibitem{AP25}
D.~J.~Aldous and B.~Pittel,
The critical beta-splitting random tree: heights and related results,
\emph{Ann. Appl. Probab.} \textbf{35} (2025), 158--195.

\bibitem{BCM25}
A.~Brandenberger, B.~Chin and E.~Mossel,
Asymptotics for the harmonic descent chain and applications to critical beta-splitting trees,
arXiv:2505.24821v3.

\bibitem{Ma26}
H.~Ma,
The height of discrete-time critical beta-splitting trees,
arXiv:2609.14329v2, 2026.

\bibitem{Son26}
H.~Song,
Asymptotic transfer, variances and limit laws for Aldous's beta-splitting trees,
preprint, 2026; verification code at \url{https://doi.org/10.5281/zenodo.23188439}.

\end{thebibliography}
\end{document}
